\begin{figure*}[t]
\centering
\figplaceholder[2.0in]{\textwidth}{图 4——\bench 的 14 个任务}{%
\textbf{版式：}通栏，两行各七个等宽子图，每个任务一格，使用 \bench 渲染的主视角画面（不用照片）。每组子图上方有一条细的任务族标签条。第 1 行：\emph{抓取}（谐波、传送带、平移、旋转、遮挡）与\emph{放置}（移动容器、旋转容器）。第 2 行：\emph{穿越}（导线穿越、完整导线穿越）、\emph{捕获}（小球、桌面小球、滑块、罐体）与\emph{撞击}（撞击后抓取）。每格下方标注任务名。
\textbf{每格：}裁剪到目标运动的区域，各格比例尺相同；叠加目标的 3--4 个半透明历史位姿与一个运动方向箭头，使运动形式可见（往复、直线、旋转、消失在遮挡物后、滚动、撞击后偏转）；障碍物描边、驱动物标注，配色与图~\ref{fig:principles}c 的角色配色一致。
\textbf{规则：}不画时间轴（图~\ref{fig:architecture} 与图~\ref{fig:principles} 已有）；各格使用同一相机。}
\caption{\textbf{\bench。}五个任务族共 14 个任务，每格叠加目标近期的运动。观测、场景运动与执行共用一个物理时钟，因此策略计算时场景继续运动。}
\label{fig:bench_overview}
\end{figure*}